% SPDX-FileCopyrightText: 2026 Hasan Melih Akbulut
% SPDX-License-Identifier: CC-BY-4.0

\chapter{Verification}
\label{ch:verif}

\section{Five layers, and the failure they all share}
\label{sec:verif-shape}

The verification programme has five layers: cocotb simulation of blocks
and of the whole system, textual and netlist \emph{guards} that read the
tree without simulating it, SymbiYosys property sets over individual
blocks, an instruction-set proof of the management core against
riscv-formal, and fault-injection campaigns at register-transfer and
gate level. They are five layers and not one because each of them can be
green about something it never looked at, and the corpus has the receipts.
\dcite{41}{section 6.6} says this repository has been bitten eight times
by a green check being read wider than what it examined, and names each
by document. (Two guard files in \code{sw/tests} cite that same section
as listing \emph{nine}; the section itself says eight, and both
statements are in the tree.)
\code{sw/tests/test\_traceability.py}'s own docstring calls its own
rewrite ``the eighth instrument in this tree found to have that shape'',
and \dcite{78}{section 3.2} records the third instance found inside this
project's own guards --- the first one found in a guard written to
prevent the shape.

Four standing rules govern how results are reported, and they are worth
stating before any number. \textbf{Every assert set ships with cover
obligations, and a target with passing asserts and failing covers is
red} (\dcite{09}{section B.1}). \textbf{A bounded result is never
reported as proven}, and every bounded result carries its depth
(\dcite{09}{C.4 item 6}). \textbf{Timing is outside formal scope}, with
OpenSTA the separate authority, so formal results hold for any
timing-clean implementation (\dcite{09}{C.4 item 2}); and \textbf{no
clock frequency is claimed for this design at all}, by
\dcite{05}{section 4 rule 5}, which \dcite{53}{section 9.4} applies to
both of the numbers a layout has reached. Periods, slacks and corner
names are reported where they are measured, in
\Cref{ch:phys}; nothing in this chapter is a frequency. Finally,
\textbf{formal proves mechanisms correct, fault injection measures
architectural sensitivity, and only beam data speaks to cross-sections}
(\dcite{09}{C.4 item 4}) --- so no rate appears anywhere below.

One more convention, inherited from \code{docs/64}: a superseded measurement
is reported as superseded, with its date, rather than overwritten. Two
counts in this chapter are superseded and both are given twice.

\section{Simulation: the cocotb suites}
\label{sec:verif-cocotb}

Counted in the working tree on 15 September 2026: \textbf{33 cocotb test
modules carrying 413 test functions}, across \code{hw/tb} (10 modules,
166 functions), \code{hw/soc/tb/cocotb} (22 modules, 242) and
\code{tt/test} (1 module, 5). They are driven by 30 makefiles, and
\code{scripts/run\_cocotb.sh} reaches all of them: 7 of the 10 under
\code{hw/tb} by default, all 19 under \code{hw/soc/tb/cocotb}, 3 more
modules that share another suite's makefile and are named in the
runner's own table, and \code{tt/test}. The declared-function count is a
floor rather than the number a run reports, because several makefiles
elaborate one module at more than one parameter point:
\code{Makefile.secded} runs the encoder and the decoder,
\code{Makefile.tmr} runs widths 1, 8 and 32, \code{Makefile.scrub} runs
three parameter triples, and \code{Makefile.soc\_pnp} runs two buses.
The last recorded whole-tree run is \code{verification-log.tsv}'s row of
2026-09-11T00:55Z at \code{587fdd8}: \textbf{462 cocotb tests passing, 0
failing, 0 suites not clean}. That row is four days old at the time of
writing and is quoted as what it is.

Three suites are skipped by default and the reason is the freeze.
\code{hw/tb/Makefile.fi} regenerates \code{hw/tb/fi\_campaign\_results.json},
which \code{docs/34} pins by git blob hash for the TTIHP26b submission; running it inside a verification pass rewrote that file's
\code{wall\_seconds} field and dirtied a frozen tree. The campaign is
reproducible on demand and its result is committed, so the runner
declines to re-derive it and says so; \code{RUN\_FROZEN=1} includes it.
The two gate-level suites are skipped for a different reason --- they
need artefacts from a hardening run rather than from the source tree.

\subsection{The runner is itself an instrument, and it has a repair history}

\code{scripts/run\_cocotb.sh} exists because \code{pytest} at the
repository root collects \code{sw/tests} only, so quoting a pytest total
as ``the tests pass'' is wider than what was measured. Four things in it
are the residue of measured defects, and they are worth reading as a
compressed history of how a test count goes wrong.

It \textbf{parses each suite's results XML rather than trusting the make
exit code}, because a cocotb suite can report all tests passing and then
segfault in Icarus teardown --- observed on \code{Makefile.soc\_clint}
and reproducible on committed RTL. That is a tool defect, and the XML is
the evidence that survives it.

It \textbf{refuses to run twice at once} rather than queueing. Counting
is by ``the results XMLs written since this marker'', which orders by
time and not by owner, so two concurrent runs each count the other's
output. On 2026-09-08 two overlapping runs recorded 558 and 624 cocotb
tests for a 401-test repository, and both landed in
\code{verification-log.tsv} looking exactly like measurements; the rows
are still there, superseded by the 22:35Z row for the same commit
\code{1ecf509} with the note
\code{supersedes=two-invalid-rows-for-1ecf509-concurrent-run-overcount},
because a log that quietly loses its own bad measurements is not a log.

The marker is a \textbf{file touched after a pause} rather than a
second-resolution timestamp, and \code{docs/66} measured what the timestamp
version cost: the previous suite's XML, written in the same second the
timestamp was taken, was counted again under the next suite's name, and
\code{soc\_qspi} reported 47 tests for a 16-test suite because
\code{soc\_npu}'s 31 had landed three seconds earlier. The runner also
sums \emph{every} XML written during a suite rather than taking the
newest, because a parameterised suite writes several and taking one made
\code{soc\_gptimer} read 10 on one run and 3 on the next while ``0
failed'' stayed true throughout.

Finally it carries a \textbf{completeness check}: every
\code{test\_*.py} in those two directories must be named by some
makefile or by the runner's own table, and one that is not is an error
at the moment it is added rather than a suite that quietly never runs.
On 2026-09-11 three modules did not run --- \code{test\_aer\_fifo},
whose makefile is named \code{Makefile} and so fell outside the
\code{Makefile.*} glob, and \code{test\_soc\_npu\_defects} and
\code{test\_soc\_wdog\_strap}, each reachable only by typing
\code{MODULE=} by hand. All three were green, all three were invisible,
and an edit re-opening the defect two of them exist to catch would have
passed every check in the tree.

\subsection{What the suites cover, and what only they can see}

The block suites drive a module's ports and check its contract:
\code{test\_soc\_bus} (16 functions), \code{test\_soc\_clint} (17),
\code{test\_soc\_qspi} (16), \code{test\_soc\_wdog} (15),
\code{test\_soc\_gpio} (14), \code{test\_soc\_uart} (14),
\code{test\_soc\_boot} (18), \code{test\_soc\_npu} (35 --- the largest
in the tree). Two of them are \emph{campaigns rather than tests}:
\code{test\_soc\_wdog\_fi} carries 220 injections inside three test
functions and \code{test\_ibex\_regfile\_secded} carries 124 inside six
(\dcite{44}{section 8.3}), which is how the mechanisms' behavioural
recovery is measured on top of the proof of the code. Two more suites
have the same shape for their own blocks: \code{test\_soc\_boot\_fi}
injects into the boot register block's own state, and
\code{test\_soc\_clint\_fi} into the architectural time base's 72 stored
bits, each establishing an unhardened baseline and the delta against it
in the same file.

Three suites exist only to fail against the RTL as it was before a fix
--- \code{test\_soc\_npu\_defects}, \code{test\_soc\_wdog\_strap},
\code{test\_soc\_uart\_defects} --- and their headers document a command
line with \code{MODULE=} on it so they can be run twice, once against
the tree and once against a copy of the pre-fix source. That is why they
share a parent makefile rather than acquiring one each: two places to
keep sources and parameters in step is one place too many.

What simulation sees that nothing below it does is \textbf{a whole
machine executing a program}. \dcite{38}{section 7} runs a self-checking
C program on the sv2v output of Ibex under Icarus with no SystemVerilog
anywhere in the flow: eleven checks covering fetch, the RV32I ALU, every
RV32M operation including the ISA's divide-by-zero results, loads and
stores at three widths, branches, recursion, compressed code, CSR access,
an illegal-instruction trap, and --- test 10 --- PMP enforcement, a
locked read-only NAPOT region permitting a load and faulting a store with
\code{mcause = 7}. All eleven pass in \num{131694} cycles. On the
hardened system the equivalent run is 22 checks, stage~1 at cycle
\num{174128}, the core asleep after \num{185443} cycles, 587 console
characters and 0 framing errors, and \dcite{43}{section 9.1} reports it
cycle-identical to the run before the register file was substituted.
That cycle-identity is the strongest available evidence that a
substitution was transparent --- and \dcite{44}{section 9} is explicit
that it is an expectation and not a measurement, which is why a campaign
was run anyway.

What simulation cannot see is everything the design is that this program
never exercises: one workload, one seed, one parameter point. The gate
level adds a second limit of its own, and it is a tool limit rather than
a design one. The IHP cell models drive their outputs from
\code{delayed\_D}, \code{delayed\_CLK} and \code{delayed\_RESET\_B} ---
the negative-timing-check outputs of \code{\$setuphold} and
\code{\$recrem} --- which Icarus 12 does not implement. Every
\code{sg13g2\_dfrbpq\_1} output then sits at \code{x} and the design
reads back zeros. On the shuttle submission's own suite that is 1 of 5
passing (\dcite{15}{section 5.4}); in \code{hw/tb} it was 21 of 21
failing (\dcite{24}{section 3.1}). Both read exactly like a dead
netlist and neither is one. \code{hw/tb/Makefile.gl} refuses to run
under Icarus 12; \code{tt/test} is frozen and generated and cannot carry
a guard of its own, so \code{run\_cocotb.sh} forces \code{GATES=} on its
command line where it beats an exported \code{GATES=yes}, and applies
the version check on its behalf under \code{TT\_GATES=1}.

\section{Guards: the checks that do not simulate}
\label{sec:verif-guards}

A \emph{guard} in this project is a pytest that reads the tree --- RTL
text, a generated netlist, a flow script, a document, the git index ---
and asserts a property of it. Guards live in \code{sw/tests}, and on
15 September 2026 pytest collects \textbf{526 tests across 25 files}
there. The last recorded whole-suite run is \code{verification-log.tsv}'s
2026-09-11 row at 509 passing, 0 failing, which the count above
supersedes; \dcite{44}{section 8.3} reports 315 at its own writing and
\code{docs/43} closed at 295, and both are left standing as records of their
date. The largest files are \code{test\_doc\_links.py} (96 tests),
\code{test\_soc\_synthesis\_guards.py} (68), \code{test\_lif\_core.py}
(41), \code{test\_synthesis\_guards.py} (40) and \code{test\_secded.py}
(37).

Run whole on 15 September 2026 the suite reports \textbf{523 passed and
3 failed in 881 seconds}, and the three failures are more informative
than the 523. \textbf{Two of them are the suite catching the tree
mid-edit.} A guard that compares the newest formal task count stated in
the index and the roadmap against a count summed from the \code{[tasks]}
sections failed during the run and passes on a re-run minutes later,
because the document it reads acquired its 2026-09-15 recount while the
run was in progress. And the guard asserting that no second copy of the
frozen accelerator source exists under \code{hw/soc} failed on
\code{hw/soc/formal/soc\_npu\_prove/src/pilot\_top.v} --- which is a
SymbiYosys scratch copy, because sby copies every file of a job's
\code{[files]} section into its own working directory. That guard was
repaired the same day, and the repair is this chapter's principle
applied once more: a directory is a solver run directory when it
contains sby's own \code{config.sby}, walked up from the file, because
the run directories are called \code{<job>\_bmc}, \code{<job>\_prove},
\code{soc\_npu}, \code{regfile\_scrub} \emph{and the next job will
invent another spelling}. \textbf{The third failure is real and still
open as this is written}: \code{docs/82} landed on 15 September and is not
yet named in \code{docs/00-index.md}, so the link guard reports it
unreachable from the entry point --- which is the guard doing exactly
what it is for.

Guards exist because, as \code{sw/tests/test\_soc\_npu\_guards.py} puts
it in its own header, the cocotb suite and the formal job check
\emph{behaviour}, and neither can notice that a constant has been copied
instead of generated, that the frozen directory has been written to, or
that the design has grown a second copy of a register map. The boot
guards' header states the placement rule the whole family follows: a
guard belongs at the stage that can satisfy it. \code{soc\_boot.sby}
proves the boot counter is not writable and says nothing about
\code{rst\_ni} on that block being the \emph{system} reset, which is what
makes the counter count boots; \code{test\_soc\_boot.py} drives the
block's ports by hand and says nothing about the ports being connected;
the whole-system run boots and says nothing about the RAM sweep being
first, because on a healthy simulation an unswept RAM is
indistinguishable from a swept one.

\subsection{The principle: check the property, not the name}

The governing discipline is \code{docs/78}'s, and it was learned twice in
that document alone.

\code{scripts/spdx\_check.py} reported itself as licence-tagged when it
carried no header at all. The first version searched the first 4,000
characters for the identifier string, anywhere, in any context --- and
the file's own docstring and its own regular expression both contain
that string, so \textbf{the checker read its own prose as its own
licence}. It went unnoticed because it passed. It surfaced only when an
unrelated edit pushed the docstring past the 4,000-character window and
the file abruptly turned up missing. The repair narrows the property: a
tag counts only on a comment line within the first twelve lines, which
is where a header is and where prose about headers is not.

The second is sharper because it is a guard failing at guarding.
\code{scripts/gen\_public\_mirror.py} had to rewrite one sentence in a
CI workflow before publishing it, and refused to run if the sentence had
changed. It was written as \code{if wf in files:}, anchored on the
\emph{filename} \code{.github/\allowbreak workflows/\allowbreak docs.yml}. When three workflows
were folded into one and that file was deleted, the guard did not fire
--- it was skipped --- and the generator wrote 511 files as though
nothing were owed. \dcite{78}{section 3.2} states the rule that came out
of it: \textbf{a check that cannot fail when the thing it checks is
absent is not a check.} It now raises on absence and anchors on a
paragraph's content rather than on a path.

The same rule decided an anchor in the hardening guards.
\code{test\_soc\_shipped\_netlist\_guards.py} identifies a
triple-modular-redundancy replica by walking the fan-in cone of
\emph{the voter's own input pin}, and \dcite{75}{section 4.3} explains
why the obvious anchor is wrong: anchoring on the bank's output wire
\code{qb} --- the same net in a correct design --- misses the defect,
because in a design where the voter reads replica A twice, \code{qb}
still exists and is still driven by replica B; it is simply not what the
voter reads. The guard tries \code{u\_prot\_vote.in\_b} first, falls back
to \code{qb} only where that name did not survive, and reports which
candidate it used on every row.

\subsection{Three guards that caught something}

\paragraph{The voter wired to one replica twice.} This is the sharpest
of them because every other instrument in the repository passes on the
defect. \code{test\_a\_voter\_wired\_to\_one\_\allowbreak replica\_twice\_\allowbreak is\_caught\_only\_by\_the\_cone}
changes one argument in a scratch copy of \code{soc\_wdog.v} ---
\code{.in\_b(qa)} where the design has \code{qb} --- and
\dcite{75}{section 4.3} measures what each instrument then says
(\Cref{tab:verif-cone}). The \code{(* keep *)} attribute on the bank's
storage keeps replica B's flip-flops in the netlist although nothing
reads them, so the mutant is a design with three intact banks whose vote
is $\mathrm{maj}(q_a, q_a, q_c) = q_a$: \textbf{an upset in replica A is
not masked at all}. Every count in this repository passes on it, every
functional test passes, every proof passes, and the cone census fails.
The shape is named precisely in that section: not an instrument that
undercounted, but an instrument that counted the right number of the
wrong thing.

\begin{table}[tbp]
\centering
\caption{The mutation that only a connectivity census can see: a voter
whose second input is wired to the first replica
(\dcite{75}{section 4.3}). Every flip-flop count agrees; the union of
the three voter-input cones does not.}
\label{tab:verif-cone}
\small
\begin{tabular}{@{}lrr@{}}
\toprule
Check & intact & mutant \\
\midrule
Flip-flops per replica, by instance path & 29 / 29 / 29 & 29 / 29 / 29 \\
Total flip-flops in the block            & 132          & 132 \\
Fan-in cone per voter input              & 29 / 29 / 29 & 29 / 29 / 29 \\
\textbf{Union of the three cones}        & \textbf{87}  & \textbf{58} \\
\textbf{Flip-flops shared by two inputs} & \textbf{0}   & \textbf{29} \\
\bottomrule
\end{tabular}
\end{table}

\paragraph{A fault line that carried a bus signal.}
\code{test\_soc\_clkgate\_guards.py} discharges an assumption that a
proof cannot discharge about itself. \code{hw/soc/formal/clkgate\_wake.sby}
proves a transformation of the accelerator's clock-gate enable from two
side conditions, one of which (its assumption A2) is that the slow half
of the enable is a function of registers that are not being clocked ---
a term that cannot pulse and vanish while the block's clock is stopped.
The guard elaborates the module, flattens it, walks the fan-in cone of
\code{npu\_act\_slow} cut at every sequential cell, and fails if any
input port is reachable. It found \code{C\_INJ\_OVF}, the one fault line
that carries \code{psel\_i}, which the file's own header records that no
reading of the source had noticed.

\paragraph{A guard whose arithmetic was wrong for two widths.}
\dcite{43}{section 9.5} records a latent defect that only a width change
could expose. A synthesis guard asserted that turning the mixing
transform off on both mixed replicas loses $2 \times \mathit{collided}$
flip-flops, where \emph{collided} is the number of bits on which one
polarity mask is zero. At a protected word of 22 bits that is
$2 \times 11 = 22$ and it was right; at 29 bits it is $2 \times 15 = 30$
and the answer is 29. The design was right and the arithmetic was wrong,
and the comment above the assertion had said the right thing all along:
on \emph{every} bit one of the three replicas must collide, because
polarity offers exactly two storage functions per bit and there are three
replicas. The expression is now $\mathit{intact} - \mathit{PROT\_W}$.

\subsection{What guards cannot do}

They read text and netlists, so they catch the ordinary edit and not the
determined one --- the supervisor guards' own header says they cannot
see a watchdog kick reached through a function pointer or an unknown
macro. They compare port \emph{lists} and not port \emph{meanings}: a
port whose semantics changed upstream without its name changing is
invisible to them and to the build, and \dcite{43}{section 3} names that
as an accepted risk. And the netlist guards read git-ignored build
products, so on a fresh clone they skip, with the command that
regenerates the evidence --- and the files say in as many words that a
skip is evidence of nothing.

One guard failure is worth recording because it is the counting layer
failing. Until 2026-09-09 \code{scripts/verify.sh} enumerated formal
results with a two-level glob over \code{hw/soc/formal/*/} and
\code{formal/*/}. riscv-formal does not put its jobs there: its job
directories are four to six levels down under \code{hw/soc/out/}, so the
glob could not reach a single one, and the log recorded
\code{formal\_pass=118, formal\_other=0} for a tree that had 35
non-passing formal jobs in it. The enumeration is now
location-independent by construction --- every directory in the tree
holding both a \code{config.sby} and a \code{logfile.txt} --- and the
verdict is read from sby's own marker file rather than grepped out of
the log, because two jobs died on a configuration error before any
\code{DONE} line was printed and the old code's empty-string arm dropped
them silently: not counted as pass, not counted as other, not counted at
all. The old rows are left standing, and read as what they measured. On
the tree as reported the same day the counts became 425 and 35, and
\code{formal-dispositions.tsv} carries all 35 with a reason and a date:
\textbf{6 FAIL, 3 ERROR, 26 STOPPED}.

\section{The SymbiYosys property sets}
\label{sec:verif-sby}

Counted in the tree on 15 September 2026: \textbf{28 property sets
carrying 137 tasks} --- 8 sets and 54 tasks under \code{formal/} for the
frozen pilot, 20 sets and 83 tasks under \code{hw/soc/formal/} for the
system fabric. \dcite{63}{section 2} states ``87 sby tasks across 16
property sets --- eight under \code{formal/} for the pilot and eight
under \code{hw/soc/formal/}''; that was correct when it was written and
is superseded by the count above, and both are given because the
document is still in the tree saying it. All 54 pilot task directories
carry a \code{PASS} marker. Of the SoC's 83 tasks, 79 have a recorded
workdir: \textbf{76 PASS, 2 TIMEOUT, 1 UNKNOWN}, and the remaining four
have no verdict --- two were stopped by hand and two were still running
when this was written. \Cref{tab:verif-sby} is the complete list.

\begin{longtable}{@{}>{\raggedright\arraybackslash\small}p{0.195\linewidth}
                   >{\raggedright\arraybackslash\small}p{0.245\linewidth}
                   >{\raggedright\arraybackslash\small}p{0.225\linewidth}
                   >{\raggedright\arraybackslash\small}p{0.235\linewidth}@{}}
\caption{Every SymbiYosys job under \code{hw/soc/formal/}, read from the
job headers and \code{hw/soc/formal/Makefile}. Depths are sby depths.
Engines are \code{smtbmc yices} unless named. Results are the verdict
recorded in each task's workdir on 15 September 2026.}
\label{tab:verif-sby}\\
\toprule
\textbf{Job} & \textbf{What it binds} & \textbf{Tasks and depths} &
\textbf{Recorded} \\
\midrule
\endfirsthead
\toprule
\textbf{Job} & \textbf{What it binds} & \textbf{Tasks and depths} &
\textbf{Recorded} \\
\midrule
\endhead
\bottomrule
\endfoot
\code{soc\_bus} & \code{soc\_bus.v}, F1--F9; the decode proved against
the generated \code{soc\_memmap.vh} & bmc 24, prove 12, cover 24 &
3 PASS \\
\code{soc\_apb\_bridge} & \code{soc\_apb\_bridge.v}, A1--A11, clauses of
AMBA 3 APB & bmc 24, prove 12, cover 24 & 3 PASS \\
\code{soc\_apb\_wb} & \code{soc\_apb\_wb.v}, I1, W1--W8, P1--P3, D1--D2,
clauses of Wishbone B3 and APB & bmc 24, prove 12, cover 24 & 3 PASS,
2026-09-15 \\
\code{soc\_clint} & \code{soc\_clint.v}, C1--C4, P1--P3, E1--E2 (the
SECDED time base) & bmc 20, prove 10, cover 20 & 3 PASS \\
\code{soc\_wdog} & \code{soc\_wdog.v} over three \code{soc\_tmr\_bank}
replicas and a voter; W1--W5, W7, W8, W9a--c & bmc 40, prove 15,
prove\_w0 15, cover 60 & 4 PASS \\
\code{soc\_wdog\_tmr} & three replicas under one voter, with a free
fault at the voter & prove / \_w4 / \_w5 / \_w21 / \_w23 / \_rst1 /
\_abc at 4, bmc 20, cover 8; \code{boolector}, \code{abc pdr} for
\_abc & 9 PASS \\
\code{soc\_busstat} & \code{soc\_busstat.v}, B1--B6 & bmc 20, prove 10,
prove\_w16 10, cover 24 & 4 PASS \\
\code{soc\_gpio} & \code{soc\_gpio.v}, P1--P7 & bmc 20, prove 8,
prove\_w16 8, cover 20 & 4 PASS \\
\code{soc\_qspi} & \code{soc\_qspi.v}, Q1--Q8 & bmc 40, prove 12,
prove\_full 12, cover 110 & 4 PASS \\
\code{soc\_scrub} & \code{soc\_scrub.v}, C1--C7 & bmc 20, prove 10,
prove\_w16 10, cover 24 & 4 PASS \\
\code{soc\_boot} & \code{soc\_boot.v}, D1--D10, at two counter widths
and at \code{HARDEN}~$=0$ & bmc 20, prove 10, prove\_full 10,
prove\_h0 10, cover 24 & 5 PASS \\
\code{soc\_mem\_ecc} & \code{soc\_mem\_ecc.v} over a four-row array,
M1--M6 & bmc 10, prove 4, cover 14 & 3 PASS \\
\code{byte\_secded} & the $(16,8)$ byte shortening of the codec, four
per RAM row & bmc / prove / prove\_abc / cover at 4; \code{abc pdr} for
prove\_abc & 4 PASS \\
\code{regfile\_secded} & the $(40,32)$ shortening the register file
stores & bmc / prove / prove\_abc / cover at 4; \code{abc pdr} for
prove\_abc & 4 PASS \\
\code{soc\_npu\_ser} & \code{soc\_npu\_ser.v}, the transport's P, H and
D properties, at two half-period settings & bmc 30, prove 16, cover 220,
prove\_h3 16 & 4 PASS \\
\code{soc\_npu} & \code{soc\_npu.v} with its whole submodule closure;
H5 over all reachable states & bmc 30, prove 16, cover 120 & 3 PASS,
2026-09-14 \\
\code{clkgate\_wake} & \code{clkgate\_wake.v}, an abstraction; T1, T2,
W1--W4, and the same three tasks with assumption A2 dropped & bmc 24,
prove 12, cover 24, each also in an \code{\_upset} group & 6 PASS \\
\code{clkgate\_wake\_gnt} & \code{clkgate\_wake\_gnt.v}, the
wakefulness-qualified grant; G1--G3, T2 & bmc 24, prove 12, cover 24 &
3 PASS, 2026-09-15 \\
\code{regfile\_scrub} & \code{ibex\_regfile\_secded.v} at
\code{SCRUB}~$=1$ with the \emph{real} codec; R1--R4 & prove 40, bmc 40,
cover 40, prove\_pdr, bmc\_btor; \code{yices}, \code{abc pdr},
\code{boolector} & cover PASS; prove\_pdr and bmc\_btor TIMEOUT at
their bounds; prove and bmc KILLED by hand after 12\,h \\
\code{regfile\_scrub\_abs} & the same wrapper with the codec abstracted
to \code{secded\_identity.v} & prove 40, bmc 40, cover 40, prove12 at
12, prove\_pdr unbounded; \code{abc pdr} for prove\_pdr & prove\_pdr
PASS (44\,s), cover PASS, prove12 UNKNOWN; prove and bmc still running \\
\end{longtable}

\subsection{The design rules the table encodes}

\textbf{Properties live beside the design, not inside it.} Every fabric
block's properties sit in \code{<block>\_props.v}, included at the end of
the module body under \code{`ifdef FORMAL}, so the Icarus and Verilator
fault-injection flows never see them. Two jobs bind from outside
instead --- \code{soc\_mem\_ecc} and the register-file scrub pair --- and
\code{regfile\_scrub\_props.v}'s header records why that is dangerous: the
first version bound R2 and R3 hierarchically from an outside wrapper and
Yosys silently declared fresh free wires, so the trace showed the scrub
enable high in a cycle the core was writing, which the RTL cannot
produce. Those two properties now live in a file included \emph{into} the
module.

\textbf{A narrow parameter point for the covers, and a second task at the
shipped width.} \code{soc\_busstat} is proved at a counter width of 4 and
again at 16, and its header gives the reason plainly: the saturation
covers are the point of the job, a counter needing 65,535 increments
cannot be driven to its top in a bounded model, and a job that could not
reach the cover would report a green result for a case it never examined.
\code{soc\_gpio} does the same at 4 and 16 bits, \code{soc\_qspi} at
narrow and shipped frame and divider widths, \code{soc\_scrub} and
\code{soc\_boot} likewise. Nothing in any property file mentions either
number; every property is written against the parameter.

\textbf{A copied parameter is guarded by reading the RTL back.}
\code{soc\_wdog\_tmr.sby} rebuilds the composition \code{soc\_wdog.v}
builds, because a fault model needs free inputs at the voter and the
block has no port for them, and three things are therefore copied: the
three polarity masks and which replicas carry the mixing transform. The
makefile's \code{wdog\_tmr\_params} target reads all six back out of
\code{hw/soc/rtl/soc\_wdog.v} and fails if any has moved, and it runs
first. \code{boot\_tmr\_params} does the same for the boot block, which
instantiates the identical composition at a 23-bit word --- which is why
\code{soc\_wdog\_tmr} is proved at 22, 21, 23, 5 and 4 bits: a width it is
not proved at is a width it does not cover.

\textbf{Engine diversity is a task, not a footnote.}
\dcite{09}{C.4 item 7} asks for a second engine family on key proofs, and
\code{prove\_abc} in three jobs is that request expressed as a task.
\code{soc\_wdog\_tmr}'s \code{prove\_abc} additionally strips the
\code{keep\_hierarchy} attribute in that task's model only --- the AIGER
path cannot carry a hierarchical cell --- and never in the RTL, because
that attribute is one of the two defences that stop Yosys merging the
three replicas.

\textbf{And the jobs say what they do not prove.}
\code{soc\_wdog\_tmr}'s header records that every property in it stays
true of a design in which the three banks were hashed into one, because
the anti-merge argument is a property of the mapped netlist and stays
with the guards. \code{soc\_mem\_ecc}'s says the important half is
missing: nothing in that harness injects a fault, because a formal
harness that deposits into a row is proving a property of the harness.
\code{byte\_secded}'s says its fault model is one error vector over
sixteen stored bits and nothing about transients in the codec's own
gates.

\subsection{Two gaps in the makefile, one deliberate}

The makefile's \code{all} target runs 18 of the 20 jobs.
\code{regfile\_scrub} is excluded on purpose and the file says why: every
engine stalls on the codec's parity network, and a target that never
returns is not a regression test. \code{clkgate\_wake\_gnt.sby} is not
named in the makefile at all --- no target, no mention --- so its three
recorded passes of 2026-09-15 are reproducible only by invoking sby on
the file directly. That is a gap and it is stated here rather than left
to be discovered.

\section{riscv-formal on the management core}
\label{sec:verif-rvf}

\dcite{09}{part B track 1} named riscv-formal in August 2026 and nothing
was run with it until \code{docs/63}. That document is the bring-up, and it
is two phases: the stock Ibex with upstream's register file, and the same
proofs against the SECDED-protected register file the system ships.

\subsection{Binding it, which is where the work was}

\dcite{09}{B.1} recorded that Ibex exposes RVFI trace ports. True of the
SystemVerilog; \code{hw/\allowbreak soc/\allowbreak gen/\allowbreak ibex\_top.v}, the Verilog-2005 this
system is actually built from, contains \textbf{zero} occurrences of the
string \code{rvfi}, because the ports are inside \code{`ifdef RVFI} and
the sv2v invocation passes only the two defines upstream's own synthesis
flow passes. So a \emph{second} tree was generated rather than the first
one changed: \code{hw/soc/genrvfi}, 34 files against 34, 16,770 lines of
Verilog-2005 against 15,873, 132 \code{rvfi} occurrences against 0, and
zero sv2v errors in both. A three-argument run of the converter into a
scratch directory diffs empty against \code{hw/soc/gen}, so every area,
timing and campaign number measured from that tree still reproduces
byte-identically.

Then the finding that decided the harness's shape: \textbf{Ibex's RVFI
block reads sixteen signals out of its submodules by hierarchical
reference}, up to three levels deep. sv2v carries
them through verbatim --- it is a translator, not an elaborator --- and
Yosys's Verilog-2005 front end treats every dotted name as an implicitly
declared net and stops, with an error saying it cannot detect the sign
and width of an auto-wire node. There is no flag for it. The route out was
already in the tree: the pinned toolchain ships a SystemVerilog front end
that resolves all sixteen when pointed at the sv2v output. The cost is
carried as an assumption --- \textbf{the front end that elaborates the
core in this proof is not the front end that elaborates it in the build}
--- and the only corroboration is \dcite{38}{section 5}, which measured
the two front ends on the same source at \SI{0.89}{\percent} of area and
two flip-flops apart, on the SystemVerilog rather than on this tree.

One construct the second front end refuses is rewritten by a tracked
script into the run's work tree rather than patched into Ibex, so that
``patches required to Ibex: zero'' survives; the script refuses to run
unless it finds the construct exactly once, so an upstream commit that
adds another is an error and not a quietly different core.

\subsection{The ledger, and the vacuity control}

Nine assumptions are ledgered in \dcite{63}{section 4.1}, and the
\code{[defines]} block writes the switchable ones into every generated
job in plain text so that a \code{grep} lists what was assumed per check
without reading any Verilog. A1 and A2 are the bus contract; A3 ties bus
errors low, and the ledger states its cost --- nothing here says anything
about the core's behaviour on a bus error, and \code{soc\_bus.v} raises
one on every unmapped address. A4 and A5 keep the PMP transparent so that
the privilege-blind instruction models can be compared at all, at the
cost that \textbf{none of this proves PMP enforcement}. A6 bounds memory
to grant in two cycles and respond in four, on the two checks that need
it. A7 substitutes a transparent clock gate, so the core keeps clocking
through the sleep a \code{wfi} enters. A8 excludes debug entry. A9 is an
adapter, not an assumption, and \S\ref{sec:verif-rvf-defects} is the
finding behind it. Interrupts are \emph{not} assumed away: all five lines
are free every cycle.

Vacuity control matters more here than anywhere else in the programme,
because every instruction check begins by assuming the instruction
retires at the check cycle. If the depth is too small for the core to get
there the assumption is unsatisfiable, there is no trace, and
\textbf{the bounded model check passes having examined nothing}. Ibex's
sequential divider takes about 37 cycles on its own, so a divide check at
the instruction default depth of 20 would be a vacuous pass. Three things
answer it. The configuration is expanded \emph{twice}, into a bmc job set
and a cover job set, and both are run. The solver's own front end reports
\code{PREUNSAT} when the assumptions at the check step are unsatisfiable,
and sby maps that to \code{ERROR} rather than to \code{PASS} --- read out
of the pinned tool's source, so the crudest form of vacuity cannot be
green on any of the 79 checks. And the cover obligation is strictly
stronger than \code{PREUNSAT}: it requires the instruction to retire at
the check cycle \emph{and not trap while doing it}, which is not a
hypothetical difference, because one early configuration retired every
instruction with the trap flag set. \textbf{79 of 79 cover jobs pass},
and the cover set cost 1,122 job-seconds against the bmc set's 7,790 ---
finding one witness is cheap and proving no counterexample exists is not.
The gap is stated rather than left to be inferred: only two of
riscv-formal's checkers carry cover statements at all, so for the other
seven consistency checks the anti-vacuity argument is \code{PREUNSAT}
alone.

\subsection{What closed, and what did not}

\begin{table}[tbp]
\centering
\caption{riscv-formal on Ibex, both phases, at the depths of
\dcite{63}{section 8.1} and under the nine assumptions of its section
4.1. Phase 2 substitutes the SECDED register file and changes nothing
else: the two generated job files differ only in which file supplies the
register module.}
\label{tab:verif-rvf}
\small
\begin{tabular}{@{}lll@{}}
\toprule
 & \textbf{Phase 1}, stock & \textbf{Phase 2}, SECDED \\
\midrule
bmc, 79 jobs   & 70 PASS, 3 FAIL, 6 stopped
               & 69 PASS, 3 FAIL, 6 stopped, 1 ERROR \\
cover, 79 jobs & 79 PASS & 79 PASS \\
Instruction checks (70) & 62 PASS, 2 FAIL, 6 stopped & identical \\
Consistency checks (9)  & 8 PASS, 1 FAIL & 7 PASS, 1 FAIL, 1 ERROR \\
Wall clock at \code{-j8} & 50\,min 52\,s & 54\,min 45\,s \\
CPU as attributed        & 9{,}076\,s    & 14{,}027\,s ($\times 1.55$) \\
Cover set, job-seconds   & 1{,}122       & 2{,}154 ($\times 1.92$) \\
\bottomrule
\end{tabular}
\end{table}

\Cref{tab:verif-rvf} is the result. What is now proved that was not
proved before: for \textbf{62 of the 70 instructions of RV32IMC}, on
every execution the core can reach from reset within that check's depth,
the instruction's effect on the architectural state --- destination
register and value, source register numbers, next PC, memory address,
byte masks, store data, and whether it traps --- agrees with
riscv-formal's model of the ISA. Add eight consistency properties over
the retirement stream, and a bounded \emph{proof} about Ibex's bus
behaviour: within every depth in the table the core never has more than
two requests in flight on either port. Before this, the evidence for any
of it was \dcite{38}{section 7}'s one self-checking program, eleven
checks, one path through the machine.

The phase-2 cell reading \code{1 ERROR} carries a correction of 2026-09-09
that is worth repeating, because it is about vocabulary. The table
originally said \code{1 TIMEOUT}. The job's own log says
\code{DONE (ERROR, rc=16)} after a \code{SIGTERM} at 1\,h 41\,m 26\,s,
its generated configuration has no timeout option, and nothing in the
flow could have sent that signal. \textbf{TIMEOUT means a bound was set,
the job reached it, and the measurement is complete; ERROR from an
outside kill means no bound was set and no measurement was completed.}
Calling it a timeout would silently promote an interrupted run into a
finished experiment. The fix in both directions is the same: read the
runner's return code before choosing the noun.

\subsection{Five defects, and the one in the specification}
\label{sec:verif-rvf-defects}

Not one of the five is in Ibex's execution of an instruction. Three are
in this project's harness: an assumption that guarded two CSR ranges and
needed five, found because a machine-mode lockdown bit made every
instruction fetch fault; a bus response permitted in the same cycle as
the grant, whose counterexample is a \emph{real} reachable deadlock ---
a zero-latency memory hangs this core, which is a fact anyone
integrating it needs; and an assertion in the harness that bounded the
data port at one outstanding request when a misaligned access makes it
two. The fourth is a genuine mismatch between what Ibex's RVFI means and
what riscv-formal's RVFI means, in two halves: the interrupt flag marks
only interrupt entries and not exception entries, and the reported next
PC is the ALU's branch target, which is the wrong answer for a return
from trap. Neither is an execution bug and neither is harmless; both are
worked around by A9, which relaxes the PC checks across a handler entry
or a return and conceals nothing else.

The fifth is the one that argues for running the tool rather than
reasoning about it, and it is in the \emph{specification}.
\code{insn\_div\_ch0} failed on \code{div x31, x6, x25} with
$-347 \div 1{,}075{,}056{,}129$: Ibex answers 0, which is correct for
signed division truncating toward zero, and riscv-formal's model answers
3, which is what 4,294,966,949 divided by 1,075,056,129 gives ---
\emph{the same bit pattern read as unsigned}. \code{insn\_rem\_ch0}
failed the same way independently. The cause is Verilog signedness
propagation in three lines of upstream's model: the operands of a
conditional operator are context-determined, the expression type is
unsigned if either is, one branch is an unsigned concatenation, and it
therefore demotes the branch carrying the signed casts. This was checked
in a second tool rather than argued from the language reference --- Icarus
Verilog, which shares no code with Yosys, evaluates upstream's expression
on the counterexample's operands to 3 and the same division on its own
self-determined wire to 0. Phase 2 then reproduced the defect
independently on fresh operands with negative \emph{divisors} where phase
1 had a negative dividend, so the diagnosis rests on four counterexamples
from two runs.

Two consequences are recorded and neither is comfortable. Correcting the
specification made the proof \emph{harder}: against the defective model
the two checks failed in 852\,s and 731\,s, because finding one
counterexample is a satisfiability question; against corrected models
they had proved nothing after 40\,min 43\,s and 48\,min 52\,s. And the
residual risk is stated plainly in \dcite{63}{section 6}: a model that is
wrong in the same direction as the core would be invisible to this entire
document, and nothing in the flow reduces that.

\subsection{The M extension, and \code{reg\_ch0}}

\textbf{All eight instructions of the M extension have no formal result
in this project.} Six --- the four multiplies, \code{divu} and
\code{remu} --- were stopped after 35 to 40 minutes each in phase 1
because they occupied six of eight job slots; \dcite{63}{section 22},
added 2026-09-15, re-ran all six on both cores under a second solver with
a 7,200\,s bound written into the job, so that the status is the tool's
own word. \textbf{None of the twelve closes.} The covers are the finding
there: all six cover jobs pass, but at 1,388 to 10,999 seconds against 9
to 73 seconds for the same covers under the first solver, so the second
bit-blasting solver is between roughly 20 and 150 times slower on the
witness problems and a change of SMT back end is not the route out. The
other two, \code{div} and \code{rem}, produced a verdict against a model
that was wrong. So over the M extension \textbf{the delta between the two
phases is undefined, not clean}: if the substituted register file changed
\code{mulh}'s result, this experiment would not know.

\code{reg\_ch0} is the check phase 2 exists for. It carries two free
constants --- which architectural register, and which instruction in the
retirement stream --- so the solver picks both, and it asserts that the
value the core reads out of a register is the value the last instruction
to write it wrote. Between that write and that read, the substituted file
encodes the word to a 40-bit codeword, holds it in flip-flops, may scrub
it, and decodes and corrects it on the way out. \Cref{tab:verif-regch0}
is what it cost.

\begin{table}[tbp]
\centering
\caption{\code{reg\_ch0} against check depth on both register files
(\dcite{63}{section 18.4}, and section 20 for the depth-25 cells). The
stock column is not monotonic because these checks assert at \emph{one}
cycle, so a deeper check is a different satisfiability problem and not a
superset of a shallower one.}
\label{tab:verif-regch0}
\small
\begin{tabular}{@{}rlll@{}}
\toprule
Check cycle & Stock register file & SECDED register file & Ratio \\
\midrule
12 & 5\,s   & 6\,s                       & 1.2 \\
15 & 14\,s  & 15\,s                      & 1.1 \\
18 & 31\,s  & \textbf{353\,s}            & \textbf{11.4} \\
21 & 66\,s  & TIMEOUT at 14{,}400\,s     & $>218$ \\
23 & 411\,s & TIMEOUT at 14{,}400\,s     & $>35$ \\
25 & 189\,s alone, 399\,s in the run & TIMEOUT at 14{,}400\,s, six engines & $>36$ \\
\bottomrule
\end{tabular}
\end{table}

The cliff is between check cycle 15 and 18. Six engines were given the
depth-25 case under a four-hour bound and every one reached it; the
control matters more than the six rows, because on the \emph{stock} file
only two of those six close the check at all, so four of the six say
nothing by failing the hard case. The two that matter both time out.
\dcite{63}{section 17.2} states the consequence in both directions:
nothing there may be read as evidence that the substitution preserves the
property at cycle 25, and nothing may be read as evidence that it breaks
it. What does close is cycles 12, 15 and 18 --- deeper than a write and a
read, and shallower than one scrub pass, which is 31 cycles when the core
is idle.

\dcite{63}{section 23}, added 2026-09-15, takes the two routes that
document itself ranked second and third, and the result is the
register-file scrub pair in \Cref{tab:verif-sby}. Route~2 asks only the
register file: \code{regfile\_scrub.sby} binds the real
\code{ibex\_regfile\_secded.v} at \code{SCRUB}~$=1$ and states R1 (a word
written and not rewritten reads back at every later cycle while the scrub
walks), R2 (the pointer visits every register 1--31 and never 0), R3 (the
scrub and the core never contend for the write port) and R4 (a fault-free
file raises no error). \textbf{Route 2 alone does not close, and the
reason is the codec.} Cover passes at depth 40 with all 31 visits
reached, and bmc and prove sit at step 4 --- the first cycle at which R1
checks a stored word --- because the solver is being asked to relate
eight parity trees at the write encoder to eight at the read decoder
through thirty-two scrub encoders on the same storage. Parity is the
textbook hard case for conflict-driven SAT and none of the pinned solvers
carries Gaussian elimination, so the stall is structural. Two further
engines with a stated bound both report \code{TIMEOUT} at 10,800\,s; the
first pair had no verdict after twelve hours and was stopped by hand,
which is \textsc{killed} in this document's vocabulary and is recorded as
such.

\textbf{Route 3 on top of route 2 closes it, unbounded.}
\code{regfile\_scrub\_abs.sby} is the same wrapper with the codec
replaced by stand-ins in \code{hw/soc/formal/secded\_identity.v} --- same
module names, same ports, a degenerate code with one check bit that says
``the word is non-zero''. Under it, \code{abc pdr} proves R1--R4 for
every reachable state in \textbf{44 seconds}, and the cover reaches all
31 visits. The composition is the usual one: a property of the wrapper
that holds for every codec satisfying $\mathrm{dec}(\mathrm{enc}(x)) = x$
holds for the real one, whose $\mathrm{dec}(\mathrm{enc}(x)) = x$ is
\code{regfile\_secded.sby}'s separate theorem. What it does not cover is
any property that depends on the code's \emph{distance} --- R4 under an
injected upset, for instance --- and none is claimed.

Two things the abstraction taught are recorded in the stand-in's header
because a future stand-in must honour them. The first version was the
true identity code and bmc failed at step 3 with every read inverted:
\code{0x80000000} written, \code{0x7fffffff} read. The cause is in the
wrapper and not the stand-in --- the fast-correction read path flips
every bit whose syndrome column matches, and with an all-zero column set
and a zero syndrome that is every bit. So the register file carries two
structural assumptions about its codec that no property had ever stated:
every column of the parity-check matrix is non-zero, and the encoding of
the all-zero word is zero, because the file resets its check bits to a
literal zero. The real code satisfies both by construction; the
degenerate one was chosen to. Note also what the recorded
\code{prove12 UNKNOWN} in \Cref{tab:verif-sby} means: k-induction at
$k=12$ fails its induction step on R1 --- the property needs an inductive
invariant the wrapper does not state --- while the unbounded IC3 proof
succeeds. Those are not contradictory results; they are two methods, and
only one of them is a proof.

The cost ledger for the whole riscv-formal exercise is
\dcite{63}{section 21}: \textbf{304,049 core-seconds, 84.5 core-hours,
300 verdicts} --- and about \SI{92}{\percent} of that solver time
produced no verdict at all, almost all of it on one check on one register
file.

\section{Fault injection}
\label{sec:verif-fi}

\begin{table}[tbp]
\centering
\caption{The fault-injection campaigns of record. The register-transfer
campaigns run each injection twice, once with the watchdog armed and
once held off; a simulation count is given only where the document
states one. The instrument measures and the classifier judges, and they
are different files --- a division kept unchanged since \code{docs/16}.}
\label{tab:verif-fi}
\small
\begin{tabularx}{\linewidth}{@{}llrl>{\raggedright\arraybackslash}X@{}}
\toprule
Campaign & Level & Injections & Simulations & What it isolates \\
\midrule
\code{docs/42} & RTL & 1{,}300 & --- & the unhardened core, 13 strata
weighted by 2{,}198 bits \\
\code{docs/43} A & RTL & 1{,}400 & 2{,}800 & the SECDED register file,
same draws, byte-identical workload image \\
\code{docs/43} AW & RTL & 1{,}400 & 2{,}800 & the watchdog's cadence
contract on top of A \\
\code{docs/44} & RTL & 1{,}400 & 2{,}800 & the fast-correction path and
the fault-report port \\
\code{docs/46} & RTL & 812 & 1{,}624 & the window, on a partitioned
supervisor workload, 58 draws per stratum \\
\code{docs/74} & gate & 1{,}424 & 1{,}739 & \code{docs/43}'s records replayed
on the sign-off netlist \\
\code{docs/75} & gate & 174 & 4 netlists & the reset on a corrected upset:
two layouts and both designs' synthesis netlists \\
\bottomrule
\end{tabularx}
\end{table}

\subsection{At register-transfer level}

The campaigns share one instrument and one classifier, and each draw is
seeded per stratum and index, which is what makes a re-run a
\emph{paired} comparison rather than a new experiment.
\dcite{43}{section 8.1} checks rather than assumes it: the workload
compiles to a byte-identical ROM image, the clean run is the same 18,682
cycles with the same signature and the same injection window, and a
twenty-line check prints \code{True} on the claim that 1,300 draws are
identical draw for draw to the earlier campaign's.

The register file is the result the hardening was bought for. With the
watchdog held off --- the core with no backstop at all --- the two
register-file strata go from 78 masked, 0 corrected, 3 silently corrupt
and 4 hung in 100 draws, to 6, 94, 0 and 0; the new stratum covering the
protection's own 253 flip-flops gives 4, 96, 0, 0. All 190 corrected
injections ended with the golden answer, and there were \textbf{zero
uncorrectable syndromes in 1,400 injections}. Design-weighted, silent
corruption falls from \SI{2.8}{\percent}~$\pm$~1.6 to
\SI{1.3}{\percent}~$\pm$~0.4 and hangs from \SI{2.1}{\percent} to
\SI{0.3}{\percent}. Two cautions travel with that table in
\dcite{43}{section 8.2} and both cut against it: the \emph{denominator}
changed, because the design has 253 more flip-flops and the new stratum
takes \SI{10.3}{\percent} of the weight at a measured rate of zero, so
part of the weighted improvement is arithmetic rather than physics; and
0 of 100 is a rate below about \SI{3.7}{\percent}, not a rate of zero.

\code{docs/44} re-ran that campaign on a design with a rebuilt register read
path and got the strongest null result the harness can express: the two
1,400-record files are \textbf{byte-identical}, same md5, 27 columns
each, and the two logs differ in one line, which is the output path. That
is stronger than ``the rates agree'', because an aggregate delta of zero
is consistent with two records moving in opposite directions and this is
not: no record moved.

\code{docs/46}'s campaign is the one that produced a design decision. The
windowed watchdog rejected a kick as too early on 9 of 812 injections,
and the seven records that did not leave the core dead were replayed
verbatim against a build identical except that the window is off. \textbf{Two of them are genuine
catches} --- wrong answers in the multiplier's state register that
nothing else in the design announced, silently corrupt without the window
and announced with it. \textbf{Two of the four spurious escalations are
the window's and two are not}; the other two escalate by ordinary expiry
at the window off, because the window forced the timeout down into a band
and a shorter timeout is more easily tripped. The arithmetic is then
honest about itself: benefit 2 of 812, \SI{0.25}{\percent}
[0.07--0.89]; cost 2 of 728, \SI{0.27}{\percent} [0.08--1.00]. The
intervals overlap almost completely, and both numbers are two draws at
one site each, so they are two site behaviours and not four independent
events. \dcite{46}{section 9.1} takes the third answer --- the mechanism
is conditional and the condition should not be taken --- and records that
the feasibility argument alone carries the recommendation, because under
the other dispatch discipline no admissible window exists at any timeout
and that is arithmetic before it is a measurement.

\subsection{At gate level}

\code{docs/74} is the first gate-level simulation of the whole system in this
project, and it replays the register-transfer campaign of record one
record at a time on the netlist of the sign-off layout ---
\code{hw/soc/pnr/runs/s71boot/final/nl/soc\_top.nl.v}, md5
\code{2e43d08f\ldots}, \num{10831575} bytes, 5,873 sequential cells and
six vendor macros. Choosing the layout netlist over the synthesis netlist
costs nothing that this campaign can measure, and that was measured
rather than assumed: the two files have identical flip-flop instance
names in the same order, and 5,823 of 5,873 identical Q-net names.

Before any comparison, the 1,400 records were replayed on the current RTL
to ask whether the RTL had moved under the design: \textbf{0 of 1,400
changed}, histogram for histogram and record for record. So whatever the
netlist campaign says differently is a difference between two levels and
not between two RTLs.

Of the 1,400 records, \textbf{1,175 have a netlist twin} the campaign can
inject through, identified by a per-cycle trace of both designs, by the
elaborated design's own alias names, or both; 225 do not. Of those 1,175,
\textbf{1,158 classify identically}, and of the 1,129 like-for-like
records --- excluding 46 injections into a flip-flop that synthesis made
the twin of two register-transfer registers at once --- \textbf{1,126 are
identical}. The 17 disagreements are two structures: 15 are injections
into the instruction register that an optimisation pass stored once for
two copies, and 2 are the multiplier's re-encoded state; the netlist is
the worse of the two every time.

The controls are where the campaign earns the right to those numbers,
and control 5 is the one to state. Four records were run at both levels
with the whole mapped flip-flop state written out per cycle and compared
over the 964 flip-flops the trace identifies: \textbf{0 of 964 differing
anywhere in 18,681 cycles}. An upset deposited into a register and an
upset forced onto its netlist flip-flop leave the two designs in the same
state, cycle for cycle, for the rest of the run. That control also caught
two things on its own: an injection-instant race that cost the first pass
of 450 records, and a controller state bit whose netlist flop carries the
register-transfer name but toggles while that bit is constant, because a
re-encoding pass kept a name. Seven records were withdrawn from the
replay by the second, and the mapping rule now refuses any alias-only
mapping that the trace contradicts.

What the redundancy did in the gates that would be manufactured: the
register file corrected 171 of 180 upsets --- the same 171, record for
record, with 0 silent corruptions and 0 uncorrectable --- and the
watchdog's three replica banks voted out 174 of 174.

\subsection{The finding only the gate level could produce}

\textbf{And in every one of those 174 the system reset itself one clock
after the upset.} \code{docs/75} is that finding. The watchdog's reset
request was a combinational decode --- five bits of the \emph{voted}
word, two of whose three replicas present an exclusive-or tree over all
29 of their flip-flops --- and the top level hung the reset synchroniser
off it asynchronously, under a comment asserting that the signal was
registered in that clock domain. It was not. The netlist says it is a
glitch three ways: the settled decode is correct on all 174 records, the
same deposit at register-transfer level does not reset, and a census of
every asynchronous-reset fan-in set in the design shows that four of the
five are driven by nothing or by registered signals while the fifth is
driven by 25 flip-flops of the three replica banks and nothing else.

The repair, W9, is one flip-flop: the same predicate computed from the
pre-storage word instead of the voted one, and-ed with the decode.
\textbf{And no proof, no simulation and no register-transfer campaign in
this repository can tell the two designs apart}, because they are the
same function of the same state on every cycle --- \code{soc\_wdog.sby}
proves exactly that as property W9a by k-induction at depth 15, which is
what makes the ``no behaviour change'' claim a theorem rather than a
comment. So the evidence had to come from two other places. On the mapped
netlist, a satisfiability check shows that no assignment to the netlist's
inputs asserts the reset request while the added flip-flop holds zero,
and two mutations that would make that vacuous both fail it. And the
campaign was re-run as a two-by-two: 174 injections into the same 87
replica flip-flops at the same 174 cycles, on the unrepaired and repaired
layouts \emph{and} on both designs' synthesis netlists, mapped on
separate occasions with no clock tree and no placement. \textbf{Resets
per corrected upset: 174 of 174 before, 0 of 174 after, on both pairs.}
The row is the design and the column is nothing.

There is also a cocotb test for W9, and its own docstring states what it
cannot do: it cannot fail on the design W9 replaced. What it catches is
W9 implemented \emph{wrongly}, and that mutation is measured --- a
flip-flop registered off the decode instead of off the pre-storage
function fails exactly one of the suite's 15 tests, and the other
fourteen pass on the mutant.

\section{The cross-check, and where the layers meet}
\label{sec:verif-cross}

\dcite{09}{target \#5} is the any-state recovery template, and its
acceptance criterion says the formal result and the force/release
campaign ``must agree'', with any disagreement a blocker at gate F3.
Nothing compared them until 2026-09-14, and the reason is structural: the
comparison sits across \code{hw/tb/} and \code{formal/} and neither pass
owned it.

\code{scripts/formal\_fi\_crosscheck.py} makes the comparison.
\code{formal/lif\_ctrl.sby}'s \code{bmc\_safe} task starts a trace from an
\emph{illegal} state encoding and proves the design reaches its safe
state --- a task its own header records is not redundant with the
induction proof, because the inductive invariant ``the encoding is always
legal'' makes the recovery antecedent unreachable in the induction step,
so a default arm that resumed in idle would pass \code{prove} and fail
\code{bmc\_safe}. The campaigns inject into exactly that encoding, so
every such record is an empirical instance of the theorem's antecedent,
and the theorem forbids one outcome class: silent corruption. Run on the
tree today, the script compares \textbf{24 injections across two arms ---
12 gate-level and 12 register-transfer, all 12 of each classified
DETECTED} --- and reports agreement. \code{sw/tests/test\_artefact\_digests.py}
is what makes it run, because a script nobody runs is not a gate; the
same file checks that every run tree the paper cites is named in the
artefact manifest, for the same reason.

The script prints what it does not do, which is half of the item: it
compares the theorem's prediction with the campaign's outcomes, and does
not check that the netlist the gate-level arm injected into is the one
the proof elaborated. Closing that is an equivalence obligation, and
\textbf{no equivalence job exists in this repository} --- the same
sentence \dcite{09}{B.2 layer 3} has to write about its own target \#2,
repeated in \dcite{63}{section 6} and listed sixth in that document's
ranked next steps.

\section{What each layer sees, and where the blind spots are}
\label{sec:verif-layers}

The argument for five layers is that each sees a class of defect the one
below it structurally cannot, and the corpus has a measured instance of
each.

\textbf{A proof sees what a simulation cannot: all reachable states.}
\code{soc\_npu.sby}'s H5 --- that the accelerator's input strobe is
exactly one state of the event engine, equal and not merely implied ---
had been checked by one cocotb test and one textual guard;
\dcite{55}{section 9.3} named the gap and declined to close it and
\dcite{56}{section 9.4} declined again, in those words, before the job
was written. \code{regfile\_scrub\_abs.sby} proves the read-after-write
contract for every register, every value and every reachable state, which
is what the core-level check could not do at depth 25 in four hours under
six engines.

\textbf{A simulation sees what a proof cannot: a whole machine running a
program.} The proofs are per block; nothing proves the composition. The
cycle-identical whole-system run, the 22 checks and the 587 console
characters are the only evidence that the blocks fit together.

\textbf{A guard sees what neither can: the netlist, the wiring and the
tree.} All twenty formal tasks stayed green on a design where the
synthesiser had collapsed three banks into one (\dcite{41}{section 9.4}),
and the only check that could see it counted flip-flops in the mapped
netlist. The voter-input cone census sees the case a flip-flop count
cannot reach at all. And a guard is the only layer that can notice that
the frozen directory has been written to, or that a constant was copied
instead of generated.

\textbf{A fault-injection campaign sees what none of them can: what an
upset actually costs this workload.} The proofs say the mechanisms are
correct; the campaign says that the class of failure a hardening was
bought for went from 3 silent corruptions and 4 hangs in 100 draws to 0
and 0.

\textbf{And the gate level sees what the register-transfer level cannot,
in both directions.} 17 records classify differently, every one of them
worse on the netlist, because an optimisation pass stored two registers
once. And the reset on a corrected upset exists only in the gates: it is
invisible to every proof, every simulation and every register-transfer
campaign in this repository, by construction rather than by oversight.

The blind spots that remain, stated as the documents state them:

\begin{itemize}
\item \textbf{No equivalence job exists.} The register-transfer design
and the mapped netlist have never been formally related, and the
front-end split of \S\ref{sec:verif-rvf} has now added a second place
where one would be worth having.
\item \textbf{No CSR is checked.} riscv-formal's CSR checks bind to ports
Ibex does not carry, so nothing says anything about the machine status,
trap vector, exception PC, cause, interrupt or PMP registers --- which is
what the intended software architecture rests on. Closing it means a
patch to Ibex.
\item \textbf{PMP enforcement is not proved.} The hardware is elaborated
and in the load/store path of everything proved above, and the evidence
that it enforces anything is one simulated test: one locked read-only
region, one permitted load, one faulting store.
\item \textbf{The M extension has no formal result}, under two solvers
and at two stated bounds.
\item \textbf{Nothing anywhere injects a single-event transient.} Every
campaign forces or deposits into storage. The register file's decoder
alone is 528 combinational cells, the voter is combinational, the fetch
address adder is combinational, and a processor is mostly not flip-flops.
\item \textbf{No timing, anywhere in this chapter.} Formal results hold
for any timing-clean implementation and the gate-level simulations use
zero-delay cell models, so a fault whose effect depends on a setup or
hold miss is invisible --- and the layout this campaign ran on does not
meet its setup constraint (\dcite{70}{section 6}). The physical glitch W9
repairs therefore has a width that is unmeasured and unmeasurable here.
\item \textbf{One workload, one seed, one parameter point}, deliberately,
so that records could be replayed. Nothing says what a program that used
the PMP or took interrupts would see.
\item \textbf{No rate.} Fault injection measures architectural
sensitivity; only beam data speaks to cross-sections.
\item \textbf{Some evidence is only as good as the tree it is run in.}
The netlist and coverage guards read git-ignored build products and skip
on a fresh clone; the shuttle submission's gate-level arm needs a netlist
nothing in this repository produces and refuses rather than falling back
to a register-transfer run.
\end{itemize}

Two of those --- the transient and the equivalence obligation --- are the
ones that would change what the rest of this chapter means, and neither
is scheduled.
